% Source: report/sections/01_introduction.md (owner: Janu, #3).
\section{Introduction}
\label{sec:intro}

Colorectal cancer is one of the most common cancers worldwide and one of the leading causes of cancer
death~\cite{sung2021globocan}. Most colorectal cancers grow out of \textbf{polyps}, small growths on the inner
wall of the colon that can turn malignant over several years. Finding and removing polyps during a
\textbf{colonoscopy}, a video examination of the colon, is therefore the main way to prevent the disease.

Colonoscopy is not perfect: a meta-analysis of tandem colonoscopy studies found that about a quarter of
adenomas are missed~\cite{zhao2019miss}, and every percentage point gained in a doctor's adenoma detection
rate is linked to a lower risk of cancer appearing before the next screening~\cite{corley2014adr}. The
polyps missed most often are the \textbf{small} ones and the \textbf{flat (sessile)} ones, which barely stand out
from the surrounding tissue.

Computer assistance can help here. \textbf{Semantic segmentation} labels every pixel of a colonoscopy frame as
polyp or background. Compared with only answering ``is there a polyp?'', a pixel mask shows the doctor the
polyp's exact location, size and outline, which supports both detection and the decision on how to remove it.
For this to be useful during a live examination, the model must also be fast enough to keep up with the
video stream, typically 25--30 frames per second.

In this project we study \textbf{U-Net}~\cite{ronneberger2015unet}, the standard architecture for biomedical
segmentation, and several of its variants on the public \textbf{Kvasir-SEG} dataset~\cite{jha2020kvasir}. Each
step changes one idea, so its effect can be measured:
\begin{enumerate}
  \item a \textbf{U-Net trained from scratch}, the baseline;
  \item a U-Net whose encoder is \textbf{pretrained on ImageNet} (ResNet-34), plus a lightweight \textbf{MobileNetV2}
        encoder, to measure the value of transfer learning;
  \item an \textbf{Attention U-Net}~\cite{oktay2018attention}, which learns where to look, and
        \textbf{SegFormer}~\cite{xie2021segformer}, which replaces the convolutional encoder with a transformer.
\end{enumerate}

\paragraph{Contributions.}
\begin{itemize}
  \item A controlled comparison of five U-Net variants on the same fixed split, with the same augmentation,
        loss and evaluation code, reported with Dice and IoU.
  \item An \textbf{accuracy versus inference-speed} study on a single T4 GPU at batch size 1, which shows which models
        can run in real time during a colonoscopy, including the cost of test-time augmentation.
  \item A \textbf{failure analysis on small and flat polyps}, using polyp size, expert sessile/flat labels from
        Kvasir-Sessile~\cite{jha2021sessile} and a colour-contrast measure, together with miss rates.
  \item A \textbf{cross-dataset test} on CVC-ClinicDB~\cite{bernal2015cvc}, images from another hospital, to show how
        well each model generalises without retraining.
  \item A public, reproducible code base in which each experiment runs as one notebook on a free GPU.
\end{itemize}
